\subsection{局部环上的 Macaulay 模}

命题 3.--- 设 A 为 Noether 局部环，M 为非零的有限生成 A-模，d 为 M 的维数。下列条件等价：

\begin{enumerate}
\def\labelenumi{(\roman{enumi})}
\item
  A-模 M 是 Macaulay 的
\item
  有 \(\operatorname{prof}(\mathbf{M}) = d\)；
\item
  有 \(\mathrm{Ext}_{\Lambda}^{i}(\kappa_{\mathrm{A}},\mathrm{M}) = 0\)（对每个整数 \(i < d\)）；
\item
  有 \(\mathrm{Ext}_{\mathrm{A}}^{i}(\mathrm{N},\mathrm{M}) = 0\)（对每个有限长度的 A-模 N 及每个整数 \(i < d\)）；
\item
  有 \(\operatorname{Ext}_A^i (\mathrm{N},\mathrm{M}) = 0\)（对每个有限生成 A-模 N 及每个整数 \(i < d - \dim (\mathrm{M}\otimes_{\mathrm{A}}\mathrm{N})\)）；
\item
  存在一个由 \(\mathfrak{m}_{\mathrm{A}}\) 中元素组成的长度为 \(d\) 的 M-正则序列。
\end{enumerate}

条件 (i) 等价于等式 prof(M) = d，亦即等价于条件 (ii)，或者再等价于不等式 prof(M) ≥ d，也就是等价于 (iii) 与 (vi) (§ 1, 第 4 节, 定理 2)。蕴含关系 (v) ⇒ (iv) 与 (iv) ⇒ (iii) 是显然的。

最后，设 M 是 Macaulay 的，并设 N 为一个有限生成 A-模。令 \(F = \text{Supp}(N)\)；由 II, § 4, 第 4 节, 命题 18，有 \(\text{Supp}(M) \cap F = \text{Supp}(M \otimes_{A} N)\) ，从而

\[
\operatorname{prof} _ {F} (M) = \operatorname{codim} (\operatorname{Supp} (M) \cap F, \operatorname{Supp} (M)) = \dim M - \dim (M \otimes_ {A} N)
\]

(第 1 节, 命题 1 的推论及第 2 节, 命题 2 的推论)。蕴含关系 (i)⇒(v) 于是由 § 1, 第 5 节, 注 1 得出。

在本节其余部分，我们称 Noether 局部环 A 上的有限生成模 M 是纯的，如果对 M 的每个伴随素理想 p，都有 \(\dim(\mathrm{A}/\mathfrak{p}) = \dim(\mathrm{M})\)。这也意味着 M 没有嵌入伴随素理想，并且 M 的支撑集的诸不可约分支都具有相同的维数。Noether 局部环上的每个 Macaulay 模都是纯的 (第 2 节, 命题 2 及其推论)。

引理 1.-- 设 A 为 Noether 局部环，M 为有限生成的纯 A-模，x 为 \(\mathfrak{m}_{\mathrm{A}}\) 的一个元素。下列条件等价：

\begin{enumerate}
\def\labelenumi{(\roman{enumi})}
\item
  有 \(\dim(M/xM)=\dim(M)-1\)；
\item
  乘以 x 的映射 \(x_{M}\) 是单射。
\end{enumerate}

不妨设 M 非零。断言 (i) 等价于 x 不属于 Supp(M) 中满足 dim(A/p) = dim(M) 的任何极小元 p (VIII, § 3, 第 2 节, 命题 3)，而断言 (ii) 等价于 x 不属于 M 的任何伴随素理想 (IV, § 1, 第 1 节, 命题 2 的推论 2)。由于 M 是纯的，其伴随素理想就是 Supp(M) 的极小元，故 (i) 与 (ii) 等价。

设 A 为 Noether 局部环，M 为非零的有限生成 A-模。回忆（VIII, § 3, 第 2 节）：称 \((x_{1},\ldots,x_{r})\)（其元素属于 \(m_{A}\)）这样的序列对 M 是割的，如果有 \(\dim(M/(x_{1}M+\ldots+x_{r}M))=\dim(M)-r\)。

命题 4. 设 A 为 Noether 局部环，M 为非零的有限生成 A-模，\((x_{1},\ldots ,x_{r})\) 为 \(\mathfrak{m}_{A}\) 中元素的一个对 M 割的序列。下列条件等价：

\begin{enumerate}
\def\labelenumi{(\roman{enumi})}
\item
  A-模 M 是 Macaulay 的
\item
  序列 \((x_{1},\ldots,x_{r})\) 是 M-正则的，且 A-模 \(\mathrm{M}/(x_{1}\mathrm{M}+\ldots+x_{r}\mathrm{M})\) 是 Macaulay 的。
\end{enumerate}

设序列 \((x_{1},\ldots,x_{r})\) 是 M-正则的。于是有

\[
\dim (\mathrm{M}) = r + \dim (\mathrm{M} / (x _ {1} \mathrm{M} + \dots + x _ {r} \mathrm{M}))
\]

\[
\operatorname{prof} (\mathrm{M}) = r + \operatorname{prof} \left(\mathrm{M} / (x _ {1} \mathrm{M} + \dots + x _ {r} \mathrm{M})\right)
\]

(§ 1, 第 4 节, 命题 7)，由此得到蕴含关系 (ii)⇒(i)。

设 A-模 M 是 Macaulay 的，对 \(r\) 作归纳来证明 (ii)。当 \(r = 0\) 时断言是显然的。若 \(r \geqslant 1\)，则由归纳假设，A-模 \(N = M/(x_1M + \ldots + x_{r-1}M)\) 是 Macaulay 的，且有 \(\dim(N/x_rN) = \dim(N) - 1\)，因为序列 \((x_1, \ldots, x_r)\) 是割的；于是位似 \((x_r)_N\) 是单射的（引理 1），且 \(N/x_rN\) 是 Macaulay 的 (\(\mathrm{no.}\) 1, 例 3)，由此得到 (ii)。

定理 1.--- 设 \(A\) 为 Noether 局部环，M 为非零的有限生成 A-模，d 为 M 的维数，\(\mathbf{x} = (x_1, \ldots, x_d)\) 为 \(\mathfrak{m}_{\mathrm{A}}\) 中元素的一个对 M 割的序列，J 为它生成的理想。下列条件等价：

\begin{enumerate}
\def\labelenumi{(\roman{enumi})}
\item
  A-模 M 是 Macaulay 的
\item
  序列 x 是 M-正则的；
\item
  序列 x 对 M 是完全割的 (A, X, 第 157 页, 定义 2)；
\item
  M 关于理想 J 的重数 (VIII, § 7, 第 1 节, 定义 1) \(e_{\mathrm{J}}(\mathrm{M})\) 等于 A-模 M/JM 的长度；
\item
  对每个满足 \(1 \leqslant i \leqslant d\) 的整数 i，A-模 \(\mathrm{M}/(x_{1}\mathrm{M}+\ldots+x_{i-1}\mathrm{M})\) 都是纯的。
\end{enumerate}

(ii) 与 (iii) 的等价性由 A, X, 第 160 页, 定理 1 的推论 1 得出。由于 A-模 M/JM 具有有限长度 (VIII, § 3, 第 2 节, 定理 1)，(iii) 与 (iv) 的等价性由 VIII, § 4, 第 3 节, 命题 4 及第 4 节, 定理 3 得出。剩下要证明 (i)、(ii) 与 (v) 的等价性。

(i)⇒(v)：若 M 是 Macaulay 的，则诸模 \(\mathrm{M}/(x_{1}\mathrm{M}+\ldots+x_{i-1}\mathrm{M})\) 中的每一个都是 Macaulay 的 (命题 4)，从而是纯的。

(v)⇒(ii)：对诸模 M/(x₁M+\ldots+\(x_{i-1}\){}M) 中的每一个应用引理 1 即得。

(ii)⇒(i)：由命题 4 得出，因为 M/JM 具有有限长度，从而是 Macaulay 的。

